\documentclass{studentnotes}

\typeout{OT7C-STARTER:STUDENTNOTES}

\setnotetitle{Starter Study Notes}
\setnotecourse{Course Name}
\setnoteauthor{Author Name}
\setnotedate{\today}

\begin{document}

\makenotetitle
\tableofcontents
\newpage

\section{Main Idea}

Use this page as a compact study handout. Start with the idea, then add the
formulae, examples, and reminders that should survive revision week.

\begin{definition}
  A model is a simplified description of a system that keeps the features needed
  for the question being answered.
\end{definition}

\begin{namedformula}{Linear relationship}
  \label{formula:linear-relationship}
  y=mx+c
\end{namedformula}

\begin{quicknote}
  Formula \formularef{formula:linear-relationship} is useful when equal changes
  in \(x\) produce equal changes in \(y\).
\end{quicknote}

\section{Worked Pattern}

\begin{example}
  If \(m=2\), \(c=1\), and \(x=3\), then \(y=2(3)+1=7\).
\end{example}

\begin{importantnote}
  Keep the statement, the formula, and the worked substitution close together.
  This makes the note useful when revising without the original textbook.
\end{importantnote}

\begin{personalnote}
  Add your own memory cue here.
\end{personalnote}

\remembernote{Check units before trusting a numerical answer.}

\end{document}
